% SEGMENT OLP-0042-S01
% Part: sets-functions-relations
% Chapter: arithmetization
% Section: From N to Z

\documentclass[../../../include/open-logic-section]{subfiles}


\begin{document}

\olfileid{sfr}{arith}{int}
\olsection{$\Nat$ இலிருந்து $\Int$ க்கு}

இரண்டு அடிப்படை உணர்தல்கள் இதோ:
	\begin{enumerate}
		\item ஒவ்வொரு முழுவையும் $n - m$ என்ற வடிவில் எழுதலாம்;
        இங்கு $n, m \in\Nat$.
		\item $n - m$ என்ற கோவையில் குறியாக்கப்பட்டுள்ள தகவலை,
        $\tuple{n, m}$ என்ற வரிசைப்படுத்தப்பட்ட ஜோடியாலும்
        அதேபோலக் குறியாக்கலாம்.
	\end{enumerate}
இயல் எண்களின் வரிசைப்படுத்தப்பட்ட ஜோடிகள் $\Nat^2$ இன்
!!{element}s என்பது நமக்கு ஏற்கெனவே தெரியும்.
$\Nat$ ஐ நாம் புரிந்துகொண்டுள்ளோம் என்றும் முன்கொள்கிறோம்.
எனவே இந்த இரண்டு உணர்தல்களின் அடிப்படையில் ஒரு முறைசாராத முன்மொழிவு:
\emph{முழுக்களை இயல் எண்களின் வரிசைப்படுத்தப்பட்ட ஜோடிகளாகக் கொள்வோம்}.

உண்மையில், இந்த முன்மொழிவு அளவுக்கு அதிகமாக முறைசாராதது.
$0- 2 = 4 - 6$ ஆக இருக்க வேண்டும் என்பது தெளிவு.
ஆனால் $\tuple{0, 2 }\neq \tuple{4, 6}$ என்பதும் தெளிவு.
எனவே $\Nat^2$ தான் முழுக்களின் கணம் என்று நேரடியாகக் கூற முடியாது.

% SEGMENT OLP-0042-S02
முந்தைய சிக்கலிலிருந்து பொதுமைப்படுத்தினால், நமக்குத் தேவையானது:
	$$a - b = c - d \text{ ஆக இருக்கும்போது, அப்போதும் அப்போது மட்டுமே, }a + d = c + b$$
(முழுக்கள் இவ்வாறுதான் \emph{செயல்பட வேண்டியவை} என்பது தெளிவாக
இருக்க வேண்டும்: இரு பக்கங்களிலும் $b$, $d$ ஆகியவற்றைக் கூட்டினால் போதும்.)
இந்தச் செயல்பாட்டை உறுதிசெய்வதற்கான எளிய வழி,
வரிசைப்படுத்தப்பட்ட ஜோடிகளுக்கு இடையில் $\Intequiv$ என்ற
சமானத் தொடர்பைப் பின்வருமாறு வரையறுப்பதாகும்:
$$\tuple{ a, b } \Intequiv \tuple{c, d}\text{ ஆக இருக்கும்போது, அப்போதும் அப்போது மட்டுமே, }a + d = c + b$$
இது ஒரு சமானத் தொடர்பு என்று இப்போது காட்ட வேண்டும்.

% SEGMENT OLP-0042-S03
\begin{prop}
$\Intequiv$ என்பது ஒரு சமானத் தொடர்பாகும்.
\end{prop}

% SEGMENT OLP-0042-S04
	\begin{proof}
	$\Intequiv$ தற்சுட்டு, சமச்சீர், கடப்பு ஆகிய பண்புகளைக் கொண்டது
    எனக் காட்ட வேண்டும்.
	
	\emph{தற்சுட்டுப் பண்பு:}
    $a + b = b + a$ என்பதால்,
    $\tuple{a, b} \Intequiv \tuple{a, b}$ என்பது தெளிவு.
	
	\emph{சமச்சீர்ப் பண்பு:}
    $\tuple{a, b} \Intequiv \tuple{c, d}$ எனக் கொள்க;
    எனவே $a + d = c + b$.
    அப்போது $c + b = a + d$;
    எனவே $\tuple{c, d} \Intequiv \tuple{a, b}$.
	
	\emph{கடப்புப் பண்பு:}
    $\tuple{a, b} \Intequiv \tuple{c, d}\Intequiv \tuple{m, n}$
    எனக் கொள்க.
    எனவே $a + d = c + b$, $c + n = m + d$.
    ஆகவே $a + d + c + n = c + b + m + d$;
    எனவே $a + n = m + b$.
    ஆதலால் $\tuple{a, b } \Intequiv \tuple{m, n}$.
\end{proof}

% SEGMENT OLP-0042-S05
இப்போது இந்தச் சமானத் தொடர்பைப் பயன்படுத்திச் சமான வகுப்புகளை
எடுத்துக்கொள்ளலாம்:

\begin{defn}
		முழுக்கள் என்பவை, இயல் எண்களின் வரிசைப்படுத்தப்பட்ட ஜோடிகளுக்கு
        $\Intequiv$ இன் கீழ் கிடைக்கும் சமான வகுப்புகள்;
        அதாவது $\Int = \equivclass{\Nat^2}{\Intequiv}$.
\end{defn}

% SEGMENT OLP-0042-S06
இந்த நிர்ணய வரையறையைப் பற்றி பல்வேறு \emph{மெய்யியல்} எதிர்வினைகள்
ஒருவருக்கு இருக்கலாம். அந்த எதிர்வினைகளைக் கருதுவதற்கு முன்,
சில தொழில்நுட்ப விவரங்களைத் தொடர்வது பயனுள்ளது.

முழுக்கள் எவை என்று கூறியபின், அவற்றின் மீதான அடிப்படைச் சார்புகளையும்
தொடர்புகளையும் வரையறுக்க வேண்டும்.
$\tuple{m, n}$ ஐ ஓர் !!a{element} ஆகக் கொண்ட,
$\Intequiv$ இன் கீழான சமான வகுப்பை
$\equivrep{m, n}{\Intequiv}$ என எழுதுவோம்.
\footnote{குறிப்பு: \olref[sfr][rel][eqv]{def:equivalenceclass} இல்
அறிமுகப்படுத்திய குறியீட்டைப் பயன்படுத்தினால், இதையே
$\equivrep{\tuple{m,n}}{\Intequiv}$ என எழுதியிருப்போம்.
ஆனால் அதைப் படிப்பது சற்றுக் கடினம்.}
அதாவது:
	$$\equivrep{m, n}{\Intequiv} = \Setabs{\tuple{a, b} \in \Nat^2}{\tuple{a, b}\Intequiv \tuple{m, n}}$$
இப்போது சில வரையறைகளை முன்வைக்கிறோம்:
	\begin{align*}
		\equivrep{a, b}{\Intequiv} + \equivrep{c, d}{\Intequiv} &= \equivrep{a + c, b + d}{\Intequiv}\\
		\equivrep{a, b}{\Intequiv} \times \equivrep{c, d}{\Intequiv} &= \equivrep{a c + b  d, a  d + b c}{\Intequiv}\\
		\equivrep{a, b}{\Intequiv} \leq \equivrep{c, d}{\Intequiv} &\text{ ஆக இருக்கும்போது, அப்போதும் அப்போது மட்டுமே, }a + d \leq b + c
	\end{align*}
(வழக்கத்திற்கேற்ப, அடிக்கோள்களை எளிதாகப் படிக்குமாறு,
`$(a \times b)$' என்பதைக் குறிக்க `$ab$' ஐப் பயன்படுத்துகிறேன்.)
இந்த வரையறைகள் \emph{செயல்பட வேண்டிய} விதத்தில் செயல்படுகின்றனவா
என்பதை இப்போது உறுதிசெய்ய வேண்டும்.
இதன் பொருளை விரித்துரைத்து முழுவதும் சரிபார்ப்பது மிகவும் உழைப்பானது;
அந்த விவரங்களை \olref[check]{sec} இல் தருகிறோம்.
ஆனால் சுருக்கமான முடிவு: அனைத்தும் செயல்படுகின்றன!

% SEGMENT OLP-0042-S07
இறுதியாக இன்னொரு செயல் மீதமுள்ளது.
இயல் எண்களைப் பயன்படுத்தி முழுக்களைக் கட்டமைத்துள்ளோம்.
ஆனால் இதனால் இயல் எண்கள் \emph{தாமே முழுக்களாக இல்லை} என்று ஆகிறது.
இதன் மெய்யியல் முக்கியத்துவத்தை \olref[ref]{sec} இல் மீண்டும் கருதுவோம்.
ஆனால் முற்றிலும் தொழில்நுட்ப நோக்கில்,
இயல் எண்களை \emph{முழுக்களாகக்} கருதுவதற்கு ஏதாவது ஒரு வழி தேவை.
கருத்து மிகவும் எளியது: ஒவ்வொரு $n \in \Nat$ க்கும்,
$n_\Int = \equivrep{n, 0}{\Intequiv}$ என நிர்ணயிக்கிறோம்.
இந்த வரையறை ஒழுங்காகச் செயல்படுகிறது என்பதை உறுதிசெய்ய வேண்டும்;
அதாவது எந்த $m, n \in \Nat$ க்கும்:
\begin{align*}
	(m + n)_\Int &= m_\Int + n_\Int\\
	(m \times n)_\Int &= m_\Int \times n_\Int\\
	m \leq n &\liff m_\Int \leq n_\Int
\end{align*}
ஆனால் இவை அனைத்தும் நேரடியானவை.
எடுத்துக்காட்டாக, இவற்றில் இரண்டாவது நிறைவேறுகிறது எனக் காட்ட,
இயல் எண்களின் செயல்பாட்டைப் பயன்படுத்திப் பின்வருமாறு காரணப்படுத்தலாம்:
%\begin{proof}
%	இயல் எண்களின் செயல்பாட்டைப் பயன்படுத்தினால், தெளிவாக:
	\begin{align*}
%		(m + n)_\Int  = \equivrep{m + n, 0}{\Intequiv} = \equivrep{m + n, 0 + 0}{\Intequiv} = \equivrep{m, 0}{\Intequiv} + \equivrep{n, 0}{\Intequiv} = m_\Int + n_\Int\\
		(m \times n)_\Int  &= \equivrep{m \times n, 0}{\Intequiv} \\
		&= \equivrep{m \times n + 0 \times 0, m \times 0 + 0 \times n}{\Intequiv} \\
		&= \equivrep{m, 0}{\Intequiv} \times \equivrep{n, 0}{\Intequiv} \\
		&= m_\Int \times n_\Int
	\end{align*}
%\end{proof}
மற்ற இரண்டு நிபந்தனைகள் நிறைவேறுகின்றன என்பதை உறுதிசெய்வதைப்
பயிற்சியாக விடுகிறோம்.

% SEGMENT OLP-0042-S08
\begin{prob}
	எந்த $m, n \in \Nat$ க்கும்,
    $(m + n)_\Int = m_\Int + n_\Int$ மற்றும்
    $m \leq n \liff m_\Int \leq n_\Int$ எனக் காட்டுக.
\end{prob}
\end{document}
